\section{Introduction}\label{sec:introduction}

Let a finite simple undirected graph have $N$ vertices and be given by its adjacency matrix $A$, a table of zeros and ones that records which vertices are joined by an edge. These notes answer three questions about it. (1) What does the entry $(A^k)_{ij}$ of the $k$-th power of $A$ represent? (2) How can one tell from $A$ alone whether the graph is connected? (3) What is a path between two vertices $n_i$ and $n_j$?

The answers are as follows. (1) The entry $(A^k)_{ij}$ is the number of walks of length $k$ from $n_i$ to $n_j$ (\Cref{thm:walks}; a walk is a sequence of vertices, each adjacent to the next, with repetition allowed); it counts walks, not paths, as \Cref{ex:board} shows on a three-vertex graph. (2) The graph is connected exactly when every entry off the diagonal of $A + A^2 + \dots + A^{N-1}$ is positive (\Cref{thm:test}). (3) A path between $n_i$ and $n_j$ is a walk from $n_i$ to $n_j$ in which no vertex occurs twice (\Cref{def:path}).

\Cref{sec:definitions} fixes the notation and the definitions, in the textbook's order (the course textbook, in Armenian \cite{petrosyan_graph_theory}); textbook pages are printed pages, one less than the PDF page. \Cref{sec:walks} proves the answer to question 1, \Cref{sec:connectivity} the answer to question 2, \Cref{sec:examples} works both on two graphs, and \Cref{sec:limits} states what the test does not give. Where the questions and the example graph come from is told in the Notes on sources. The textbook's own word ``path'' has a different meaning from the one used here; \Cref{rem:trap} says which.

\medskip
\noindent Story for the plain-words paragraphs. A village has houses, and two-way lanes join some pairs of different houses, with at most one lane per pair. Houses that share a lane are neighbours, and a neighbour chart has a yes-or-no box for each pair of houses. A stroll goes from house to house along lanes, may pass a house again, and has its steps counted; a no-return stroll never visits a house twice. The village is joined-up when a no-return stroll links every pair of houses.

\noindent In the dictionary of this story, a house is a vertex, a lane an edge, neighbours are adjacent vertices, the neighbour chart is the adjacency matrix, a box is an entry, a stroll is a walk, its steps are its length, a no-return stroll is a path, and joined-up means connected. Before each definition, theorem, lemma, proposition, and example, a plain-words paragraph retells the statement in this story; the formal text stands without these paragraphs.
